% chirality.tex — Section 3: Part B — Cl(6) Chirality and the Standard Model
% ASSERT_CONVENTION: natural_units=dimensionless, jordan_product=(1/2)(ab+ba), clifford_convention=euclidean_positive, octonion_basis=fano_e1e2=e4, complex_structure=u_equals_e7, spin_representation=S10plus_boyle

%----------------------------------------------------------------------
\section{Part B: Cl(6) Chirality and the Standard Model}\label{sec:chirality}
%----------------------------------------------------------------------

The second half of the derivation shows that the same algebraic
structure---the exceptional Jordan algebra probed by a C*-observer---
yields the Standard Model gauge group with the correct chiral
representation.  The key mechanism is a Clifford subalgebra $\Cl{6}$
induced by a single additional choice: a unit imaginary octonion
$u\in S^6$.

%----------------------------------------------------------------------
\subsection{The octonion splitting \texorpdfstring{$\mathbb{O} = \mathbb{C} \oplus \mathbb{C}^3$}{O = C + C3}}\label{sec:oct-split}
%----------------------------------------------------------------------

Choose a unit imaginary octonion $u\in S^6 \subset \mathrm{Im}(\mathbb{O})$.
Following the Fano convention, we set $u = e_7$.

\begin{remark}[Gap~B, step~2]\label{rem:gap-B2}
The choice of $u$ is a second symmetry-breaking input, independent
of the idempotent choice $E_{11}$ (Remark~\ref{rem:gap-B1}).
All unit imaginary octonions are conjugate under
$\Gfour = \mathrm{Aut}(\mathbb{O})$; the orbit is
$S^6 = \Gfour/\SU{3}$, $\dim(S^6) = 14 - 8 = 6$.
What selects a particular $u$ is not addressed in this work.
\end{remark}

The choice of $u$ embeds $\mathbb{C}$ into $\mathbb{O}$:
\begin{equation}\label{eq:C-embed}
  \mathbb{C}
  \;=\; \mathrm{span}_\mathbb{R}\{1, e_7\}
  \;\subset\; \mathbb{O}.
\end{equation}
The orthogonal complement of $u$ in $\mathrm{Im}(\mathbb{O})$
is 6-dimensional:
\begin{equation}\label{eq:W-def}
  W \;=\; u^\perp \cap \mathrm{Im}(\mathbb{O})
  \;=\; \mathrm{span}_\mathbb{R}\{e_1,e_2,e_3,e_4,e_5,e_6\},
  \qquad \dim(W) = 6.
\end{equation}

The space $W$ inherits a complex structure from left multiplication
by~$u$:
\begin{equation}\label{eq:J-def}
  J\colon W \to W,
  \qquad
  J(w) = e_7 \cdot w.
\end{equation}
To verify $J^2 = -\mathrm{Id}$: for any $w\in W$, the left
alternative identity $x\cdot(x\cdot y) = (x\cdot x)\cdot y$ gives
\begin{equation}\label{eq:J-squared}
  J(J(w))
  \;=\; e_7\cdot(e_7\cdot w)
  \;=\; (e_7\cdot e_7)\cdot w
  \;=\; -w,
\end{equation}
since $e_7^2 = -1$.  The complex structure pairs the six real
directions into three complex coordinates:
$(e_k,\, e_7\cdot e_k)$ for $k=1,2,3$.  Consequently
\begin{equation}\label{eq:oct-split}
  \mathbb{O}
  \;=\; \mathbb{C} \oplus \mathbb{C}^3
  \qquad
  (\text{as real vector spaces: } 2 + 6 = 8),
\end{equation}
where $W\cong \mathbb{C}^3$ as a complex vector space (with $J$
playing the role of multiplication by~$i$).

\medskip
\noindent\textbf{Stabilizer.}
The stabilizer of $u$ in $\Gfour = \mathrm{Aut}(\mathbb{O})$
is precisely the group that preserves the complex structure $J$
on $W\cong\mathbb{C}^3$:
\begin{equation}\label{eq:SU3C-stab}
  \SU{3}_C
  \;:=\; \mathrm{Stab}_{\Gfour}(u)
  \;=\; \mathrm{Stab}_{\Gfour}(e_7),
  \qquad \dim(\SU{3}_C) = 8.
\end{equation}
This $\SU{3}_C$ acts on $W\cong\mathbb{C}^3$ via the defining
3-dimensional representation and will be identified with the
color group.

%----------------------------------------------------------------------
\subsection{\texorpdfstring{$\Cl{6}$}{Cl(6)} inside \texorpdfstring{$\Cl{10}$}{Cl(10)}}\label{sec:cl6}
%----------------------------------------------------------------------

The 6 directions of $W$ become generators of a Clifford subalgebra
inside $\Cl{10}$.  Recall from Section~\ref{sec:spin-upgrade} that
$\Spin{10}$ acts on $\Stenplus$ via the Clifford algebra $\Cl{10}$
with generators $\Gamma_1,\ldots,\Gamma_{10}$ satisfying
$\{\Gamma_A,\Gamma_B\} = 2\delta_{AB}$.
The 6 generators corresponding to the $W$ directions define
\begin{equation}\label{eq:cl6-gens}
  \gamma_k := \Gamma_k, \quad k=1,\ldots,6,
  \qquad
  \{\gamma_i,\gamma_j\} = 2\delta_{ij},
\end{equation}
generating the subalgebra
\begin{equation}\label{eq:cl6-def}
  \Cl{6}
  \;=\; \mathrm{Alg}(\gamma_1,\ldots,\gamma_6)
  \;\subset\; \Cl{10}.
\end{equation}

The remaining 4 generators $\Gamma_7,\Gamma_8,\Gamma_9,\Gamma_{10}$
correspond to the complement: the $u = e_7$ direction plus three
additional directions from the $\Vzero$ Peirce space
(the $h_2(\mathbb{O})$ sector).

\emph{The $\Cl{6}$ subalgebra is not introduced ad hoc}: it is
induced by the choice of complex structure $u$, which selects the
6 directions $e_1,\ldots,e_6\in W = u^\perp\cap\mathrm{Im}(\mathbb{O})$.

\medskip
\noindent\textbf{Volume form.}
The volume element of $\Cl{6}$ is
\begin{equation}\label{eq:omega6}
  \omegasix \;=\; \gamma_1\gamma_2\gamma_3\gamma_4\gamma_5\gamma_6.
\end{equation}

\begin{proposition}[Properties of $\omegasix$]\label{prop:omega6}
\begin{enumerate}
  \item[\textup{(a)}] $\omegasix^2 = -1$.
  \item[\textup{(b)}] The chirality operator $i\,\omegasix$ satisfies
    $(i\,\omegasix)^2 = +1$.
  \item[\textup{(c)}] $P = \tfrac{1}{2}(1 - i\,\omegasix)$ is an
    idempotent projector with $\Tr(P) = 16$ on the 32-dimensional
    Dirac spinor $\Delta_{10}$ of $\Cl{10}$.
\end{enumerate}
\end{proposition}

\begin{proof}
(a)~To compute $\omegasix^2 = (\gamma_1\cdots\gamma_6)^2$, we
move the second copy of each generator past the first.
The number of transpositions is $5+4+3+2+1+0 = 15 = 6\cdot 5/2$.
Each transposition costs $(-1)$ (from $\{\gamma_i,\gamma_j\}=0$
for $i\ne j$), and each $\gamma_k^2 = 1$, giving
\begin{equation}\label{eq:omega6-sq}
  \omegasix^2
  \;=\; (-1)^{15}\cdot 1
  \;=\; -1.
\end{equation}
Equivalently, the general formula $\omega_n^2 = (-1)^{n(n-1)/2}$
gives $(-1)^{15} = -1$ for $n=6$.

(b)~$(i\,\omegasix)^2 = i^2\cdot\omegasix^2 = (-1)(-1) = +1$.

(c)~Since $(i\,\omegasix)^2 = 1$, the element
$P = \tfrac{1}{2}(1-i\,\omegasix)$ is idempotent:
$P^2 = \tfrac{1}{4}(1-2i\,\omegasix + 1) = P$.
For the trace on the 32-dimensional Dirac spinor
$\Delta_{10} = \Stenplus\oplus S_{10}^-$:
the volume form $\omegasix$ anticommutes with each $\gamma_k$
(since $\gamma_k\omegasix = (-1)^5\omegasix\gamma_k = -\omegasix\gamma_k$),
so its eigenvalues come in $\pm$ pairs, giving $\Tr(\omegasix)=0$.
Hence $\Tr(P) = \tfrac{1}{2}(32 - 0) = 16$.
\end{proof}

The projector $P$ selects a 16-dimensional subspace of
$\Delta_{10}$ corresponding to one generation of Standard
Model fermions~\cite{Todorov2022}.

%----------------------------------------------------------------------
\subsection{Pati-Salam breaking}\label{sec:pati-salam}
%----------------------------------------------------------------------

The volume form $\omegasix$ breaks $\Spin{10}$ to the subgroup
that commutes with it.  The 10 Clifford generators split into:
\begin{itemize}
  \item \textbf{6 internal generators} $\gamma_1,\ldots,\gamma_6$
    (the $W$ directions from $\Cl{6}$);
  \item \textbf{4 external generators} $\Gamma_7,\Gamma_8,
    \Gamma_9,\Gamma_{10}$ (the complement).
\end{itemize}

The $\binom{10}{2} = 45$ bivectors $\Gamma_A\Gamma_B$ ($A<B$)
spanning $\mathfrak{spin}(10)$ split into three classes based on
their commutation with~$\omegasix$:

\medskip
\noindent\textbf{(i)~Internal--internal:}
$\gamma_i\gamma_j$ for $i,j\in\{1,\ldots,6\}$, $i<j$.
Each factor anticommutes with $\omegasix$, so
$(\gamma_i\gamma_j)\,\omegasix = \omegasix\,(\gamma_i\gamma_j)$:
\emph{these commute with~$\omegasix$}.
Count: $\binom{6}{2} = 15$.
They generate $\mathfrak{spin}(6)\cong\mathfrak{su}(4)$.

\medskip
\noindent\textbf{(ii)~External--external:}
$\Gamma_A\Gamma_B$ for $A,B\in\{7,\ldots,10\}$, $A<B$.
Each external generator commutes with $\omegasix$ (passing
through 6 internal generators incurs $(-1)^6 = +1$), so
$(\Gamma_A\Gamma_B)\,\omegasix = \omegasix\,(\Gamma_A\Gamma_B)$:
\emph{these commute with~$\omegasix$}.
Count: $\binom{4}{2} = 6$.
They generate $\mathfrak{spin}(4)\cong\mathfrak{su}(2)_L\oplus\mathfrak{su}(2)_R$.

\medskip
\noindent\textbf{(iii)~Mixed:}
$\gamma_i\Gamma_A$ for $i\in\{1,\ldots,6\}$, $A\in\{7,\ldots,10\}$.
The internal factor anticommutes with $\omegasix$ while the external
factor commutes, so
$(\gamma_i\Gamma_A)\,\omegasix = -\omegasix\,(\gamma_i\Gamma_A)$:
\emph{these anticommute with~$\omegasix$} and are not in the stabilizer.
Count: $6\times 4 = 24$.

\medskip
\noindent\textbf{Stabilizer dimension check:}
$15 + 6 = 21$ (stabilizer) $+$ $24$ (complement) $= 45 = \dim\,\mathfrak{spin}(10)$.\ \checkmark

\begin{proposition}[Pati-Salam stabilizer]\label{prop:pati-salam}
The stabilizer of $\omegasix$ in $\Spin{10}$ is the Pati-Salam group:
\begin{equation}\label{eq:PS-stab}
  \mathrm{Stab}_{\Spin{10}}(\omegasix)
  \;=\;
  \frac{\Spin{6}\times\Spin{4}}{\mathbb{Z}_2}
  \;\cong\;
  \frac{\SU{4}\times\SU{2}_L\times\SU{2}_R}{\mathbb{Z}_2},
  \qquad \dim = 21.
\end{equation}
\end{proposition}

\noindent\textbf{Representation decomposition.}
Under Pati-Salam, the 16-dimensional Weyl spinor $\Stenplus$
(established in Section~\ref{sec:spin-upgrade} as the complexified
Peirce half-space $\Vhalf^{\,\mathbb{C}}$) decomposes as
\begin{equation}\label{eq:16-PS}
  \mathbf{16}
  \;\to\;
  (\mathbf{4},\mathbf{2},\mathbf{1})
  \;\oplus\;
  (\bar{\mathbf{4}},\mathbf{1},\mathbf{2}).
\end{equation}

\noindent\textbf{Dimension check:}
$4\times 2\times 1 + 4\times 1\times 2 = 8 + 8 = 16$.\ \checkmark

\medskip
The physical content of each factor:
\begin{itemize}
  \item $(\mathbf{4},\mathbf{2},\mathbf{1})$:
    4 colors ($3$ quark + $1$ lepton) $\times$ $\SU{2}_L$ doublet $\times$
    $\SU{2}_R$ singlet = \textbf{left-handed fermions}
    ($(u_L,d_L)$ in 3 colors $+$ $(\nu_L,e_L)$).
  \item $(\bar{\mathbf{4}},\mathbf{1},\mathbf{2})$:
    $\bar{4}$ colors $\times$ $\SU{2}_L$ singlet $\times$ $\SU{2}_R$ doublet
    = \textbf{right-handed fermions}
    ($(u_R,d_R)$ in 3 colors $+$ $(\nu_R,e_R)$).
\end{itemize}

This is the \textbf{LEFT (chiral) embedding}: $\SU{2}_L$ acts
\emph{only} on the $(\mathbf{4},\mathbf{2},\mathbf{1})$ component,
i.e.\ on left-handed fermions.  Right-handed fermions are $\SU{2}_L$
singlets.  This is precisely the chirality pattern of the Standard
Model.

\begin{remark}[Left vs.\ diagonal embedding~\cite{BaezSawin}]\label{rem:left-vs-diag}
In the diagonal embedding, a single $\SU{2}$ would act on
\emph{both} chiralities symmetrically, producing a non-chiral
theory.  The theorem of Sawin (as cited in~\cite{BaezSawin})
establishes that the $\Spin{10}$ GUT embedding of the SM is the
LEFT embedding, not the diagonal.  Our construction reproduces
this: the volume form $\omegasix$ separates left and right sectors
\emph{before} the Pati-Salam breaking, ensuring $\SU{2}_L$
acts only on the left-handed component.
\end{remark}

%----------------------------------------------------------------------
\subsection{SM gauge group from Pati-Salam}\label{sec:SM-gauge}
%----------------------------------------------------------------------

The same complex structure $u=e_7$ that produced $\Cl{6}$ further
breaks $\SU{4}$.  The fundamental representation $\mathbf{4}$ of
$\SU{4} = \Spin{6}$ acts on $\mathbb{C}^4$.  The complex structure
$u$ distinguishes the lepton direction (associated with $u$ itself)
from the three quark/color directions:
\begin{equation}\label{eq:SU4-break}
  \mathbb{C}^4 = \mathbb{C}_{\text{lepton}} \oplus \mathbb{C}^3_{\text{color}}.
\end{equation}
The stabilizer of this decomposition in $\SU{4}$ is
$\SU{3}\times\mathrm{U}(1)\subset\SU{4}$,
giving the Standard Model gauge group:
\begin{equation}\label{eq:SM-gauge}
  \SMgauge
\end{equation}
with hypercharge $Y = (B-L) + 2J_3^R$, where $B-L$ comes from
$\SU{4}\to\SU{3}\times\mathrm{U}(1)_{B-L}$
and $J_3^R$ from $\SU{2}_R$.

\noindent\textbf{Dimension check:}
$\dim(\SU{3}_C) + \dim(\SU{2}_L) + \dim(\mathrm{U}(1)_Y)
= 8 + 3 + 1 = 12$.\ \checkmark

\noindent\textbf{Full breaking chain:}
\begin{equation}\label{eq:breaking-chain}
  \Spin{10}
  \;\xrightarrow{\;\omegasix\;}
  \frac{\SU{4}\times\SU{2}_L\times\SU{2}_R}{\mathbb{Z}_2}
  \;\xrightarrow{\;u\;}
  \SU{3}_C\times\SU{2}_L\times\mathrm{U}(1)_Y\,.
\end{equation}
Both breaking steps are driven by the same input: the choice of
complex structure $u = e_7$ on $\mathrm{Im}(\mathbb{O})$.

%----------------------------------------------------------------------
\subsection{Quantum numbers of one generation}\label{sec:quantum-numbers}
%----------------------------------------------------------------------

Following Furey~\cite{Furey2018}, the Witt decomposition of $\Cl{6}$
defines creation and annihilation operators
\begin{equation}\label{eq:witt}
  a_j = \tfrac{1}{2}(\gamma_{2j-1} + i\gamma_{2j}),
  \qquad
  a_j^\dagger = \tfrac{1}{2}(\gamma_{2j-1} - i\gamma_{2j}),
  \qquad j=1,2,3,
\end{equation}
satisfying the canonical anticommutation relations
$\{a_i,a_j^\dagger\} = \delta_{ij}$, $\{a_i,a_j\}=0$.
The particle-number operator $N = \sum_j a_j^\dagger a_j$ has
eigenvalues $0,1,2,3$ on the $2^3 = 8$ Fock states of the $\Cl{6}$
spinor.  The volume form acts as
\begin{equation}\label{eq:omega6-witt}
  \omegasix = -i\cdot(-1)^N,
\end{equation}
so even-$N$ states have eigenvalue $-i$ and odd-$N$ states have
eigenvalue $+i$.  The projector $P = \tfrac{1}{2}(1-i\,\omegasix)$
selects one chirality sector.

The Schwinger boson construction from $\Cl{4}$ Witt operators
$b_1,b_2$ gives the Cartan generators of $\SU{2}_L\times\SU{2}_R$:
$J_3^L = \tfrac{1}{2}(m_1+m_2-1)$, $J_3^R = \tfrac{1}{2}(m_1-m_2)$,
where $m_k = b_k^\dagger b_k$.

The complete set of quantum number operators is:
\begin{align}
  B-L &= 1 - \tfrac{2}{3}N,
  &  Y &= (B-L) + 2J_3^R,
  &  Q &= J_3^L + \tfrac{Y}{2},
  \label{eq:qn-ops} \\
  T_3^c &= \tfrac{1}{2}(n_1 - n_2),
  &  T_8^c &= \tfrac{1}{2\sqrt{3}}(n_1 + n_2 - 2n_3).
  &
  \label{eq:color-ops}
\end{align}

All 16 SM states for one generation, with quantum numbers computed as
explicit matrix eigenvalues~\cite{Furey2018,Todorov2022}, are:

\begin{center}
\small
\renewcommand{\arraystretch}{1.15}
\begin{tabular}{@{}clccrrrr@{}}
  \toprule
  & Particle & $N$ & $J_3^L$ & $J_3^R$ & $B\!-\!L$ & $Y$ & $Q$ \\
  \midrule
  \multirow{4}{*}{\rotatebox{90}{$(\mathbf{4},\mathbf{2},\mathbf{1})$}}
  & $u_L$ ($r,g,b$) & 1 & $+\tfrac{1}{2}$ & $0$ &
      $+\tfrac{1}{3}$ & $+\tfrac{1}{3}$ & $+\tfrac{2}{3}$ \\
  & $d_L$ ($r,g,b$) & 1 & $-\tfrac{1}{2}$ & $0$ &
      $+\tfrac{1}{3}$ & $+\tfrac{1}{3}$ & $-\tfrac{1}{3}$ \\
  & $\nu_L$         & 3 & $+\tfrac{1}{2}$ & $0$ &
      $-1$           & $-1$             & $0$ \\
  & $e_L$           & 3 & $-\tfrac{1}{2}$ & $0$ &
      $-1$           & $-1$             & $-1$ \\
  \midrule
  \multirow{4}{*}{\rotatebox{90}{$(\bar{\mathbf{4}},\mathbf{1},\mathbf{2})$}}
  & $u_R$ ($r,g,b$) & 1 & $0$ & $+\tfrac{1}{2}$ &
      $+\tfrac{1}{3}$ & $+\tfrac{4}{3}$ & $+\tfrac{2}{3}$ \\
  & $d_R$ ($r,g,b$) & 1 & $0$ & $-\tfrac{1}{2}$ &
      $+\tfrac{1}{3}$ & $-\tfrac{2}{3}$ & $-\tfrac{1}{3}$ \\
  & $\nu_R$         & 3 & $0$ & $+\tfrac{1}{2}$ &
      $-1$           & $0$             & $0$ \\
  & $e_R$           & 3 & $0$ & $-\tfrac{1}{2}$ &
      $-1$           & $-2$            & $-1$ \\
  \bottomrule
\end{tabular}
\end{center}

\noindent\textbf{Consistency checks:}
8 $\SU{2}_L$ doublets ($J_3^L = \pm\tfrac{1}{2}$) $+$ 8 singlets ($J_3^L=0$);
4 color singlets (leptons) $+$ 12 color triplets (quarks: 3 colors $\times$ 4);
$Q$ distribution $\{+\tfrac{2}{3}\!:6,\; -\tfrac{1}{3}\!:6,\; 0\!:2,\; -1\!:2\}$
matches the Pati-Salam SM assignment;
$Y = (B-L) + 2J_3^R$ and $Q = J_3^L + Y/2$ verified for all 16 states.
Particle content: $u_L\times 3$, $d_L\times 3$, $\nu_L$, $e_L$,
$u_R\times 3$, $d_R\times 3$, $\nu_R$, $e_R$ --- complete
one-generation SM fermion content.
